% Shared source for the student sheet and the solutions sheet.
% Content inside a solution environment appears only in the solutions PDF.

\studentonly{%
  \begin{instructions}
    Give a clear justification for each answer. You may use results proved in
    the lecture unless an exercise asks you to establish them. Questions marked
    \textbf{Extension} are optional.
  \end{instructions}}

\solutiononly{%
  \begin{instructions}
    These notes reproduce each exercise before its solution. Equivalent
    arguments and notation should receive full credit when the mathematical
    reasoning is correct.
  \end{instructions}}

\begin{exercise}{From statistics to partitions}
  Let $u_1,\ldots,u_6\in\mathcal X_+^*$ be six possible histories. Suppose
  that they induce the following conditional distributions over the entire
  future:
  \[
    \begin{array}{c|cccccc}
      \text{history} & u_1 & u_2 & u_3 & u_4 & u_5 & u_6\\
      \hline
      \text{future distribution} & P & P & P & Q & Q & S
    \end{array}
  \]
  where $P$, $Q$, and $S$ are pairwise distinct.

  Three statistics $R_A$, $R_B$, and $R_C$ induce the partitions
  \[
    \begin{aligned}
      \Pi_A
        &=\bigl\{\{u_1,u_2\},\{u_3\},\{u_4\},\{u_5\},\{u_6\}\bigr\},\\
      \Pi_B
        &=\bigl\{\{u_1,u_2,u_3\},\{u_4,u_5\},\{u_6\}\bigr\},\\
      \Pi_C
        &=\bigl\{\{u_1,u_2,u_3,u_4\},\{u_5\},\{u_6\}\bigr\}.
    \end{aligned}
  \]
  Recall that a statistic is sufficient for prediction when histories assigned
  the same value induce the same distribution over the future.

  \begin{parts}
    \item What does it mean for two histories to belong to the same cell of the
      partition induced by a statistic? Illustrate your answer using $u_1$ and
      $u_2$ in $\Pi_A$.

    \item Determine which of $R_A$, $R_B$, and $R_C$ are sufficient. If a
      statistic is not sufficient, give two histories witnessing the failure.

    \item Which cells of $\Pi_A$ can be merged without losing sufficiency?
      Perform all such merges and identify the resulting partition.

    \item Explain why no two cells of $\Pi_B$ can be merged while preserving
      sufficiency. What does this show about $R_B$?

    \item Prove that every refinement of a sufficient partition is also
      sufficient. Is every coarsening of a sufficient partition necessarily
      sufficient? Use $\Pi_A$, $\Pi_B$, and $\Pi_C$ to illustrate your answer.

    \item \textbf{Extension.} Consider the partition
      \[
        \Pi_D
          =\bigl\{\{u_1\},\{u_2,u_3\},\{u_4,u_5\},\{u_6\}\bigr\}.
      \]
      Show that $\Pi_A$ and $\Pi_D$ are both sufficient but that neither
      refines the other. How can both nevertheless refine the same minimal
      sufficient partition?
  \end{parts}

  \begin{solution}
    \begin{parts}
      \item Two histories belong to the same cell when they receive
        the same summary value. In particular,
        $R_A(u_1)=R_A(u_2)$.

      \item Both $R_A$ and $R_B$ are sufficient. Every cell of $\Pi_A$ and
        every cell of $\Pi_B$ contains histories carrying only one of the
        distributions $P$, $Q$, or $S$.

        The statistic $R_C$ is not sufficient. Its first cell contains $u_1$
        and $u_4$, for example, but $u_1$ induces $P$ while $u_4$ induces $Q$.
        Since $P\neq Q$, these histories cannot be assigned the same value by
        a sufficient statistic.

      \item The cell $\{u_1,u_2\}$ can be merged with $\{u_3\}$ because all
        three histories induce $P$. Similarly, $\{u_4\}$ can be merged with
        $\{u_5\}$ because both histories induce $Q$. No cell can be merged
        with $\{u_6\}$, and no $P$-cell can be merged with a $Q$-cell. After
        all permissible merges, the resulting partition is
        \[
          \bigl\{\{u_1,u_2,u_3\},\{u_4,u_5\},\{u_6\}\bigr\}=\Pi_B.
        \]

      \item The three cells of $\Pi_B$ correspond respectively to the three
        distinct future distributions $P$, $Q$, and $S$. Merging any two cells
        would therefore place histories with different future distributions
        in the same cell and destroy sufficiency. Hence $R_B$ is minimal
        sufficient.

      \item Let $\Pi'$ refine a sufficient partition $\Pi$. Any two histories
        in the same cell of $\Pi'$ also belong to the same cell of $\Pi$.
        Since $\Pi$ is sufficient, those histories induce the same future
        distribution. Therefore $\Pi'$ is sufficient.

        In the example, $\Pi_A$ refines $\Pi_B$, and both are sufficient. A
        coarsening need not remain sufficient: $\Pi_C$ is obtained from the
        sufficient partition $\Pi_A$ by merging its cells
        $\{u_1,u_2\}$, $\{u_3\}$, and $\{u_4\}$. The resulting cell contains
        histories inducing both $P$ and $Q$, so $\Pi_C$ is not sufficient.

      \item Every cell of $\Pi_D$ lies within one predictive class, so $\Pi_D$
        is sufficient. The same was established for $\Pi_A$.

        The cell $\{u_1,u_2\}$ of $\Pi_A$ is split between $\{u_1\}$ and
        $\{u_2,u_3\}$ in $\Pi_D$, so $\Pi_A$ does not refine $\Pi_D$.
        Conversely, $\{u_2,u_3\}$ in $\Pi_D$ is split between
        $\{u_1,u_2\}$ and $\{u_3\}$ in $\Pi_A$, so $\Pi_D$ does not refine
        $\Pi_A$. Both partitions nevertheless refine $\Pi_B$, because every
        one of their cells is contained in one of the three predictive classes
        represented by $P$, $Q$, and $S$.
    \end{parts}
  \end{solution}
\end{exercise}

\begin{exercise}{Causal states are minimal sufficient statistics}
  For $x,x'\in\mathcal X_+^*$, define the relation $\sim$ by
  \[
    x\sim x'
    \quad\Longleftrightarrow\quad
    \Pr\!\left(
      X_{|x|+1:|x|+|y|}=y\mid X_{1:|x|}=x
    \right)
    =
    \Pr\!\left(
      X_{|x'|+1:|x'|+|y|}=y\mid X_{1:|x'|}=x'
    \right)
  \]
  for every $y\in\mathcal X^*$. Define the causal state of $x$ by
  \[
    \epsilon(x):=[x]_{\sim},
  \]
  and write $\mathcal S:=\epsilon(\mathcal X_+^*)$ for the set of
  causal states.

  \begin{parts}
    \item Prove that $\sim$ is an equivalence relation on
      $\mathcal X_+^*$ by establishing each of the following:
      \begin{subparts}
        \item reflexivity: $x\sim x$ for every $x\in\mathcal X_+^*$;
        \item symmetry: if $x\sim x'$, then $x'\sim x$;
        \item transitivity: if $x\sim x'$ and $x'\sim x''$, then
          $x\sim x''$.
      \end{subparts}

    \item Prove that $\epsilon$ is sufficient for prediction.

    \item Let $R:\mathcal X_+^*\to\mathcal R$ be any sufficient statistic.
      Prove that
      \[
        R(x)=R(x')\quad\Longrightarrow\quad
        \epsilon(x)=\epsilon(x').
      \]
      Deduce that every sufficient partition refines the causal-state
      partition.

    \item Using parts (b) and (c), conclude that $\epsilon$ is a minimal
      sufficient statistic.

    \item \textbf{Extension.} Let $T$ be another minimal sufficient
      statistic. Prove that $T$ and $\epsilon$ induce the same partition of
      $\mathcal X_+^*$. Deduce that there is a bijection
      \[
        \psi:\mathcal S\to\operatorname{im}(T)
      \]
      satisfying $T=\psi\circ\epsilon$.
  \end{parts}

  \begin{solution}
    \begin{parts}
      \item
        \begin{subparts}
          \item For every $x\in\mathcal X_+^*$ and every continuation $y$,
            the conditional probability after $x$ equals itself. Hence
            $x\sim x$, so $\sim$ is reflexive.

          \item If $x\sim x'$, then the two continuation probabilities are
            equal for every $y$. Reversing the equality shows that
            $x'\sim x$, so $\sim$ is symmetric.

          \item Suppose $x\sim x'$ and $x'\sim x''$. For every $y$, the
            continuation probability after $x$ equals that after $x'$, which
            equals that after $x''$. Thus $x\sim x''$, so $\sim$ is
            transitive.
        \end{subparts}
        Therefore $\sim$ is an equivalence relation.

      \item If $\epsilon(x)=\epsilon(x')$, then $x$ and $x'$ belong to the
        same equivalence class, so $x\sim x'$. By the definition of $\sim$,
        the two histories induce the same probability for every finite
        continuation $y$. This is precisely the definition of sufficiency for
        prediction, so $\epsilon$ is sufficient.

      \item Suppose $R(x)=R(x')$. Since $R$ is sufficient, $x$ and $x'$
        induce the same probability for every finite continuation. Hence
        $x\sim x'$, and equivalent histories have the same equivalence class:
        $\epsilon(x)=\epsilon(x')$.

        Thus any two histories in the same cell induced by $R$ also lie in the
        same causal state. Every cell induced by $R$ is therefore contained in
        a causal state, so the partition induced by $R$ refines the
        causal-state partition. Since $R$ was arbitrary, every sufficient
        partition has this property.

      \item Part (b) shows that $\epsilon$ is sufficient. Part (c) shows that
        its partition is refined by the partition of every sufficient
        statistic. Hence the causal-state partition is the coarsest sufficient
        partition, so $\epsilon$ is a minimal sufficient statistic.

      \item Since $T$ is sufficient, part (c), with $R=T$, shows that the
        partition induced by $T$ refines the causal-state partition. Since
        $T$ is minimal and $\epsilon$ is sufficient, the causal-state
        partition must also refine the partition induced by $T$. The two
        partitions are therefore equal.

        Define $\psi(\epsilon(x)):=T(x)$. Equality of the two partitions shows
        that this definition is independent of the representative $x$ and
        that distinct causal states receive distinct $T$-values. It also
        reaches every value in $\operatorname{im}(T)$. Hence $\psi$ is a
        bijection.
    \end{parts}
  \end{solution}
\end{exercise}

\begin{exercise}{From finite generators to causal states}
  In this exercise, $S_t$ denotes the hidden state immediately before the
  transition that emits $X_{t+1}$. All parameters lie strictly between $0$
  and $1$.

  \begin{parts}
    \item \textbf{The Zero--One--Random Process.}
      Let $S_0$ be uniformly distributed over
      $\{s_0,s_1,s_R\}$. The process has transitions
      \[
        s_0\xrightarrow{\,0:1\,}s_1,
        \qquad
        s_1\xrightarrow{\,1:1\,}s_R,
        \qquad
        s_R
        \begin{cases}
          \xrightarrow{\,0:1-p\,}s_0,\\[-1mm]
          \xrightarrow{\,1:p\,}s_0.
        \end{cases}
      \]
      An edge label $x:q$ means that the process emits $x$ and follows the
      edge with probability $q$.

      \begin{subparts}
        \item List all possible histories of length two. For each history,
          determine whether it uniquely identifies the current hidden phase.
          When it does, identify that phase.

        \item For each ambiguous length-two history, determine which
          additional observations resolve the ambiguity. Identify the
          resulting phase and explain why it remains known after every
          subsequent observation.

        \item Assume that $\epsilon(\emptyset)$, $\epsilon(0)$,
          $\epsilon(1)$, and $\epsilon(10)$ are pairwise distinct. Show that
          the causal states have shortest-length representatives
          \[
            \emptyset,\quad 0,\quad 1,\quad 10,\quad 00,\quad 01,\quad 11.
          \]
          Explain why no additional causal states arise.

        \item Under the assumption in part (iii), draw the automaton
          associated with the causal-state dynamics.
      \end{subparts}

    \item \textbf{The Even Process.}
      Let $S_0$ be uniformly distributed over $\{A,B\}$ and consider the
      generator
      \[
        A\xrightarrow{\,0:q\,}A,
        \qquad
        A\xrightarrow{\,1:1-q\,}B,
        \qquad
        B\xrightarrow{\,1:1\,}A.
      \]
      Thus every run of $1$s between successive $0$s has even length.

      \begin{subparts}
        \item For each of the length-one histories $0$ and $1$, determine
          whether the current hidden state can be identified exactly.

        \item Starting from each ambiguous history in part (i), consider its
          possible one-symbol extensions. Which extensions resolve the
          ambiguity, and which remain ambiguous? You may use the fact that
          $\epsilon(\emptyset)=\epsilon(11)$. What does this imply for
          histories of the form $1^n$?

        \item Assume additionally that
          $\epsilon(\emptyset)\neq\epsilon(1)$. Show that the causal states
          have shortest representatives
          \[
            \emptyset,\qquad 1,\qquad 0,\qquad 01,
          \]
          and explain why there are no additional causal states. You may
          distinguish causal states using possible and impossible future
          continuations.

        \item Using $\epsilon(\emptyset)=\epsilon(11)$ and assuming
          $\epsilon(\emptyset)\neq\epsilon(1)$, draw the automaton associated
          with the causal-state dynamics.
      \end{subparts}

    \item \textbf{The Simple Nonunifilar Source.}
      Let $S_0$ be uniformly distributed over $\{A,B\}$ and consider
      \[
        A\xrightarrow{\,0:1/2\,}A,
        \qquad
        A\xrightarrow{\,0:1/2\,}B,
        \qquad
        B\xrightarrow{\,0:1/2\,}B,
        \qquad
        B\xrightarrow{\,1:1/2\,}A.
      \]

      \begin{subparts}
        \item Show that observing a $1$ identifies the new hidden state as
          $A$.

        \item Starting from $A$, show that there are $n+1$ hidden paths that
          emit $0^n$, each with probability $2^{-n}$. How many of these paths
          can subsequently emit $1$? Hence calculate
          \[
            \Pr(X_{n+2}=1\mid X_{1:n+1}=10^n).
          \]

        \item Show that this probability is different for every $n\geq0$.
          Conclude that the histories
          \[
            1,10,100,\ldots
          \]
          belong to distinct causal states and that the process has infinitely
          many causal states.
      \end{subparts}
  \end{parts}

  \begin{solution}
    \begin{parts}
      \item
        \begin{subparts}
          \item All four length-two histories are possible. Their resulting
            hidden phases are
            \[
              \begin{array}{c|cccc}
                x_{1:2} & 00 & 01 & 10 & 11\\
                \hline
                S_2 & s_1 & s_R & \{s_0,s_1\} & s_0.
              \end{array}
            \]
            Thus $00$, $01$, and $11$ identify the phase exactly, while $10$
            is ambiguous.

          \item From the phases compatible with $10$, only $s_0$ can emit
            $0$, and only $s_1$ can emit $1$. Consequently, $100$ identifies
            the new phase as $s_1$, while $101$ identifies it as $s_R$. From
            each phase, the emitted symbol determines the next phase, so every
            later observation preserves synchronization.

          \item The histories $00$, $01$, and $11$ identify $s_1$, $s_R$,
            and $s_0$, respectively. These phases induce distinct next-token
            distributions: $s_1$ emits only $1$, $s_0$ emits only $0$, and
            $s_R$ can emit either symbol. Hence they give three distinct causal
            states.

            They are also distinct from the four states assumed in the
            question. This can be seen without calculating probabilities by
            comparing possible length-two continuations:
            \[
              \begin{array}{c|c}
                \text{representative} & \text{possible length-two continuations}\\
                \hline
                \emptyset & \{00,01,10,11\}\\
                0          & \{01,10,11\}\\
                1          & \{00,01,10\}\\
                10         & \{01,10,11\}\\
                00         & \{10,11\}\\
                01         & \{00,10\}\\
                11         & \{01\}
              \end{array}
            \]
            The only repeated continuation set belongs to $0$ and $10$, whose
            causal states are assumed to be distinct. Finally, every history
            of length at least three identifies the phase by part (ii), and so
            belongs to one of the states represented by $00$, $01$, and $11$.
            No further causal states arise.

          \item Appending each possible symbol gives the causal-state update:
            \[
              \begin{array}{c|cc}
                \text{state} & 0 & 1\\
                \hline
                \epsilon(\emptyset) & \epsilon(0)  & \epsilon(1)\\
                \epsilon(0)         & \epsilon(00) & \epsilon(01)\\
                \epsilon(1)         & \epsilon(10) & \epsilon(11)\\
                \epsilon(10)        & \epsilon(00) & \epsilon(01)\\
                \epsilon(00)        & \text{--}    & \epsilon(01)\\
                \epsilon(01)        & \epsilon(11) & \epsilon(11)\\
                \epsilon(11)        & \epsilon(00) & \text{--}
              \end{array}
            \]
            Here a dash denotes an impossible extension. Drawing one node for
            each row and the indicated symbol-labelled edges gives the
            required automaton.
        \end{subparts}

      \item
        \begin{subparts}
          \item The history $0$ identifies the current state as $A$: only $A$
            can emit $0$, and this transition returns to $A$. The history $1$
            is ambiguous. If $S_0=A$, it leads to $B$, while if $S_0=B$, it
            leads to $A$.

          \item From the two states compatible with the history $1$, only $A$
            can emit $0$. Therefore $10$ resolves the ambiguity and identifies
            the new state as $A$. The extension $11$ remains ambiguous. Using
            the given equality and recursively updating after pairs of $1$s
            gives
            \[
              \epsilon(1^{2k})=\epsilon(\emptyset),
              \qquad
              \epsilon(1^{2k+1})=\epsilon(1),
              \qquad k\geq0.
            \]

          \item If a history contains a $0$, its final $0$ synchronizes the
            generator to $A$. An even number of trailing $1$s then leaves it
            in $A$, represented by $0$, while an odd number leaves it in $B$,
            represented by $01$. If the history contains no $0$, part (ii)
            shows that its causal state is represented by $\emptyset$ or $1$,
            according to the parity of its length. Thus no additional causal
            states arise.

            These four states are distinct. The state represented by $01$
            cannot emit $0$, whereas the other three can. The continuation
            $10$ is impossible after $0$ but possible after $\emptyset$ and
            $1$. Finally, $\epsilon(\emptyset)\neq\epsilon(1)$ by assumption.
            Hence the displayed representatives are also shortest.

          \item Appending each possible symbol gives the causal-state update:
            \[
              \begin{array}{c|cc}
                \text{state} & 0 & 1\\
                \hline
                \epsilon(\emptyset) & \epsilon(0) & \epsilon(1)\\
                \epsilon(1)         & \epsilon(0) & \epsilon(\emptyset)\\
                \epsilon(0)         & \epsilon(0) & \epsilon(01)\\
                \epsilon(01)        & \text{--}   & \epsilon(0)
              \end{array}
            \]
            Here a dash denotes an impossible extension. Drawing one node for
            each row and the indicated symbol-labelled edges gives the
            required automaton.
        \end{subparts}

      \item
        \begin{subparts}
          \item The only transition that emits $1$ starts from $B$, and that
            transition leads to $A$. Thus the hidden state immediately after
            observing $1$ is $A$.

          \item One path remains in $A$ while emitting all $n$ zeros. Each of
            the other $n$ paths moves from $A$ to $B$ on one of the $n$
            transitions and then remains in $B$. Thus there are $n+1$ paths,
            each with probability $2^{-n}$, and
            \[
              \Pr(0^n\mid S=A)=(n+1)2^{-n}.
            \]
            The $n$ paths that end in $B$ can subsequently emit $1$. Including
            this final transition, each has probability $2^{-(n+1)}$, so
            \[
              \Pr(X_{n+2}=1\mid X_{1:n+1}=10^n)
                =\frac{n2^{-(n+1)}}{(n+1)2^{-n}}
                =\frac{n}{2(n+1)}.
            \]

          \item Since
            \[
              \frac{n}{2(n+1)}=\frac12-\frac{1}{2(n+1)},
            \]
            this probability is strictly increasing with $n$. Hence every pair
            of histories $10^n$ and $10^m$, with $n\neq m$, already has a
            different next-token distribution. They must lie in different
            causal states. The source consequently has infinitely many causal
            states despite its two-state generator.
        \end{subparts}
    \end{parts}
  \end{solution}
\end{exercise}

\studentonly{\clearpage}
\begin{exercise}{Why next-token prediction is enough}
  Recall that
  \[
    \mathcal X_+^*
      :=\left\{x\in\mathcal X^*:
        \Pr\!\left(X_{1:|x|}=x\right)>0\right\}
  \]
  is the set of allowable histories.

  Consider a model with parameters $\theta$ satisfying
  \[
    h_n=f_\theta(h_{n-1},x_n),
    \qquad
    y_n=g_\theta(h_{n-1},x_n).
  \]
  Fix the initial hidden state $h_0$ and suppose that the model predicts the
  next-token distribution exactly: for every $n\geq1$,
  \[
    y_n
      =\Pr\!\left(
        X_{n+1}\mid X_{1:n}=x_{1:n}
      \right),
      \qquad \forall\,x_{1:n}\in\mathcal X_+^*.
  \]
  For every nonempty allowable history, define
  \[
    R(x_{1:n}):=(h_{n-1},x_n).
  \]

  For $a\in\mathcal X$, define the deterministic update
  \[
    \Phi_a(h,x):=\bigl(f_\theta(h,x),a\bigr).
  \]
  For a word $z_{1:j}$, write
  \[
    \Phi_{z_{1:j}}
      :=\Phi_{z_j}\circ\cdots\circ\Phi_{z_1},
    \qquad
    \Phi_{z_{1:0}}:=\operatorname{id}.
  \]

  \begin{parts}
    \item Show that, whenever $x_{1:n}z_{1:j}$ is allowable,
      \[
        R(x_{1:n}z_{1:j})
          =\Phi_{z_{1:j}}\!\left(R(x_{1:n})\right).
      \]
      Hence show that, if $R(x_{1:n})=R(x'_{1:m})$ and both extended
      histories are allowable, then
      \[
        R(x_{1:n}z_{1:j})=R(x'_{1:m}z_{1:j}).
      \]

    \item Let $R(x_{1:n})=R(x'_{1:m})$ and fix
      $z_{1:k}\in\mathcal X^k$. As the symbols of $z_{1:k}$ are read from
      left to right, use part (a) and exact next-token prediction to show that
      for each $0\leq j<k$ such that the prefixes through $z_{1:j}$ are
      allowable,
      \[
        \begin{aligned}
        &\Pr\!\left(
          X_{n+j+1}=z_{j+1}
          \mid X_{1:n+j}=x_{1:n}z_{1:j}
        \right)\\
        &\qquad=
        \Pr\!\left(
          X_{m+j+1}=z_{j+1}
          \mid X_{1:m+j}=x'_{1:m}z_{1:j}
        \right).
        \end{aligned}
      \]
      Deduce that either the same first zero-probability extension is
      encountered after both histories, or every prefix of $z_{1:k}$ is
      allowable after both histories.

    \item In the case where every prefix is allowable, use the conditional
      chain rule to write
      \[
        \Pr\!\left(
          X_{n+1:n+k}=z_{1:k}\mid X_{1:n}=x_{1:n}
        \right)
        =
        \prod_{j=0}^{k-1}
        \Pr\!\left(
          X_{n+j+1}=z_{j+1}
          \mid X_{1:n+j}=x_{1:n}z_{1:j}
        \right).
      \]
      Write the analogous product after $x'_{1:m}$ and compare its factors.
      Deal separately with the zero-probability case from part (b), and
      conclude directly that
      \[
        \begin{gathered}
          R(x_{1:n})=R(x'_{1:m})
          \quad\Longrightarrow\\[1mm]
          \Pr\!\left(
            X_{n+1:n+k}=z_{1:k}\mid X_{1:n}=x_{1:n}
          \right)
          =
          \Pr\!\left(
            X_{m+1:m+k}=z_{1:k}\mid X_{1:m}=x'_{1:m}
          \right).
        \end{gathered}
      \]
      Hence $R(x_{1:n})=(h_{n-1},x_n)$ is a sufficient statistic for
      predicting the entire future.
  \end{parts}

  \begin{solution}
    \begin{parts}
      \item Appending one symbol $a$ gives
        \[
          R(x_{1:n}a)
            =\bigl(f_\theta(h_{n-1},x_n),a\bigr)
            =\Phi_a\!\left(R(x_{1:n})\right).
        \]
        Repeatedly applying this deterministic update along $z_{1:j}$ gives
        \[
          R(x_{1:n}z_{1:j})
            =\Phi_{z_j}\circ\cdots\circ\Phi_{z_1}
              \!\left(R(x_{1:n})\right)
            =\Phi_{z_{1:j}}\!\left(R(x_{1:n})\right).
        \]
        If $R(x_{1:n})=R(x'_{1:m})$ and both extended histories are
        allowable, applying the same composite function to the equal starting
        representations gives
        \[
          R(x_{1:n}z_{1:j})=R(x'_{1:m}z_{1:j}).
        \]

      \item Suppose that the prefixes through $z_{1:j}$ are allowable after
        both histories. Part (a) gives
        \[
          R(x_{1:n}z_{1:j})
            =R(x'_{1:m}z_{1:j}).
        \]
        The two representations give the same arguments to $g_\theta$, so
        exact next-token prediction gives
        \[
          \begin{aligned}
          &\Pr\!\left(
            X_{n+j+1}=z_{j+1}
            \mid X_{1:n+j}=x_{1:n}z_{1:j}
          \right)\\
          &\qquad=
          \Pr\!\left(
            X_{m+j+1}=z_{j+1}
            \mid X_{1:m+j}=x'_{1:m}z_{1:j}
          \right).
          \end{aligned}
        \]
        Thus the next extension has positive probability after one history
        exactly when it has positive probability after the other. Following
        the word from left to right, either both encounter their first zero at
        the same position or all its prefixes are allowable after both.

      \item If every prefix is allowable, the conditional chain rule gives
        the product in the question and
        \[
          \begin{aligned}
          &\Pr\!\left(
            X_{m+1:m+k}=z_{1:k}\mid X_{1:m}=x'_{1:m}
          \right)\\
          &\qquad=
          \prod_{j=0}^{k-1}
          \Pr\!\left(
            X_{m+j+1}=z_{j+1}
            \mid X_{1:m+j}=x'_{1:m}z_{1:j}
          \right).
          \end{aligned}
        \]
        Part (b) shows that the factors in the two products are equal, so the
        products are equal. If a first zero-probability extension is
        encountered instead, part (b) shows that it occurs after both
        histories; consequently both probabilities of the complete
        continuation are zero.

        Therefore equal values of $R$ imply equal conditional probabilities
        for every finite continuation. By the definition of sufficiency,
        \[
          R(x_{1:n})=(h_{n-1},x_n)
        \]
        is a sufficient statistic for predicting the entire future.
    \end{parts}
  \end{solution}
\end{exercise}
